\begin{theorem}[Local regular parent constraints force trivial holonomy]%
    \label{thm:peps_regular_region_parent_flatness}
    \lean{TNLean.PEPS.regionGroundSpace_regularTwisted_eq_of_internalGaugeFlat,
        TNLean.PEPS.regionPhysicalMap_mem_regularProjectorTwistedRange,
        TNLean.PEPS.regionPhysicalMap_leftInverse_on_regularGroundSpace,
        TNLean.PEPS.regularProjectorOpenRegionRange_apply_eq_zero_of_cycle_ne_one,
        TNLean.PEPS.regionPhysicalMap_apply_eq_zero_of_mem_regularGroundSpace_of_cycle_ne_one,
        TNLean.PEPS.regionPhysicalMap_slice_apply_eq_zero_of_parent_annihilates,
        TNLean.PEPS.disjoint_regularProjectorRegion_ranges_of_cycleResidual_ne_one,
        TNLean.PEPS.disjoint_regionGroundSpace_regularTwisted_of_cycleResidual_ne_one,
        TNLean.PEPS.regularWalkHolonomy_eq_one_of_mem_regionGroundSpace_of_mem_twisted,
        TNLean.PEPS.regularWalkHolonomy_eq_one_of_regionParent_annihilates_twistedState}
    \leanok
    \uses{def:peps_region_ground_space, def:peps_region_parent_interaction}
    Consider a finite graph with site-dependent regular $G$-injective
    tensors. Let $R$ induce a connected graph, and insert regular
    bond operators $u_e$ at their ordered heads. If an internal
    closed walk has nonidentity holonomy, the original regional
    ground space and the inserted regional ground space have zero
    intersection. More generally, a nonidentity spanning-tree cycle
    residual already forces this separation.

    Every regional ground vector, exposed by a local $G$-injective
    inverse, vanishes at half-edge coordinates whose cycle labels
    are not all the identity. In particular, this holds for every
    complementary slice of an arbitrary vector satisfying an
    original regional parent constraint. No inserted-bond
    representation of that vector is assumed.

    Consequently, if a nonzero inserted global PEPS vector is
    annihilated by an original regional parent term on $R$, every
    internal closed-walk holonomy in $R$ is the identity.
    If a vertex gauge removes all internal inserted operators,
    the original and inserted regional ground spaces coincide:
    crossing operators only relabel the boundary condition.
    This last assertion requires only coefficient invariance.
    No $G$-isometry is assumed. This is a necessary local implication
    in the closure argument of \cite[Theorem~5.5]{Schuch2010PEPS};
    it does not assert that every parent-kernel vector admits an
    inserted-bond description.
    This necessary local holonomy condition does not give a global closure representation; see
    \cite{gap:rmp_peps_quantum_double_g_isometry}.
\end{theorem}

\begin{proof}\leanok
    Apply a local $G$-injective inverse at every vertex of $R$.
    This exposes the canonical regular half-edge coordinates, and
    applying the original local coefficient maps restores each
    regional ground vector. In spanning-tree coordinates, the
    uninserted contraction is supported on identity cycle labels.
    The inserted contraction is supported on simultaneous
    conjugates of its cycle residuals. A nonidentity residual makes
    these supports disjoint. Since the local coordinate map has a
    left inverse on the regional ground spaces, their physical
    intersection is also zero.

    A common nonzero regional vector therefore forces all residuals
    to be the identity. Evaluating the actual closed-walk word at
    these residuals gives identity holonomy. Finally choose a
    nonzero physical slice of the inserted global vector. It belongs
    to the inserted regional ground space by contraction, and to
    the original one by the original parent constraint.
\end{proof}
